\documentclass[10pt,a4paper]{amsart}
\usepackage[margin=19mm]{geometry}
\usepackage{amsmath,amssymb,mathtools}
\usepackage[scr=rsfs,cal=euler]{mathalfa}
\usepackage{xfrac}
\usepackage[unicode,hidelinks,bookmarks=false]{hyperref}
\newcommand{\ie}{i.e.\ }
\newcommand{\cf}{cf.\ }
\newcommand{\resp}{respectively\ }
\newcommand{\Subsection}{\subsection}
\providecommand{\nobreakdash}{\nobreak\hbox{-}}
\setlength{\emergencystretch}{2em}
\multlinegap0pt
\usepackage{enumerate}
\usepackage{amssymb}
\usepackage{tikz}
\newtheorem{thm}{Theorem}
\title{Engel and Co-Engel graphs of finite groups}
\author{Peter J. Cameron, Rishabh Chakraborty, Rajat Kanti Nath, Deiborlang Nongsiang}
\date{Source: arXiv:2408.03879v5 (CC BY 4.0)}
\begin{document}
\maketitle

\noindent\emph{Source and licence.} Engel and Co-Engel graphs of finite groups. Version arXiv:2408.03879v5 (CC BY 4.0), distributed by arXiv under the Creative Commons Attribution 4.0 International licence, http://creativecommons.org/licenses/by/4.0/, as recorded in the arXiv licence metadata for this exact version and shown on the abstract page https://arxiv.org/abs/2408.03879v5. Full licence notice: \texttt{licenses/arxiv-2408.03879-CC-BY-4.0.txt}.\par\smallskip

\noindent\textit{Editorial scope.} Counted unit: the article's thm thm (source0.tex lines 596--607). The article's complete original proof of it is delivered with it (source0.tex lines 609--684, one \texttt{proof} environment of 76 lines). No mathematics is written, repaired or re-derived: the statement and the proof are the article's own lines, in the article's own notation, and no retained line is substituted or re-flowed.  No further source-local context is needed: every cross-reference of every retained line resolves inside the two delivered spans.  Nothing was omitted from any retained span, so the package carries no omission record.  No retained line contains a citation, so the wrapper carries no bibliography and no bibliography entry.  No retained line contains a figure environment, an image-inclusion command or drawing code, so no image is delivered and the package declares no asset dependency.  Editorial formatting only, in addition: an AMS \texttt{amsart} preamble that restates the article's own declarations and macros used by the retained lines, namely the environments \texttt{thm} on separate counters, together with the standard packages those lines need.  The retained environments print in this document as Theorem 1 in order of appearance, whereas the article prints the same items with the numbering of its own full document.  The complete native source of the article as distributed in the arXiv e-print for this exact version is bundled unchanged as \texttt{sources/163483/source0.tex}.
\begin{thm}
Let $D$ be a finite digraph, whose arcs are labelled with positive integers.
Then the following two conditions are equivalent:
\begin{enumerate}
\item If an arc $x\to y$ has label $1$, then there is also an arc $y\to x$
with label $1$.
\item There is a finite group $G$ and an embedding of the vertex set of
$D$ into $G$ with the properties that $x\to y$ is an arc with label $k$ if
and only if $(x,y)$ has Engel depth $k$, and there is no arc $x\to y$ if and
only if $(x,y)$ has infinite Engel depth.
\end{enumerate}
\end{thm}

\begin{proof}
First we show that (b) implies (a). The pair $(x,y)$ has Engel
depth $1$ if and only if $x$ commutes with $y$; this is equivalent to saying
that $y$ commutes with $x$, so also $(y,x)$ has Engel depth~$1$.

\medskip

To prove the converse, assume that (a) holds; for every possible $2$-vertex
digraph, we are going to construct a group $G$ with distinguished elements
$x$ and $y$ with the properties that there is an arc $x\to y$ with label $k$
if and only if $(x,y)$ has Engel depth $k$ (and no arc if and only if the
Engel depth is infinite), and similarly with $x$ and $y$ interchanged.

First we make an observation about direct products. Suppose that
$x_i,y_i\in G_i$ for $i=1,2$. Then the Engel depth of $(x_1x_2,y_1y_2)$
in $G_1\times G_2$ is the maximum of the Engel depths of $(x_1,y_1)$ in $G_1$
and $(x_2,y_2)\in G_2$. The proof is clear.

Consider the dihedral group 
\[G=\langle a,b\mid a^{2^n}=1,b^2=1,b^{-1}ab=a^{-1}\rangle\]
of order $2^{n+1}$. Now $[b,a]\in A$, so $[b,a,a]=1$; this $(a,b)$ has
Engel depth $2$. On the other hand, $[a,b]=a^{-2}$, whose order is
$2^{n-1}$; and each further commutation with $b$ halves the order, so the
Engel depth of $(b,a)$ is $n$.

Let $b'=ab$. Then $[b,b']=a^{-2}$, with order $2^{n-1}$; as above, we find
that $(b',b)$ has Engel depth $n$. The same holds for $(b,b')$.

So we can realise the case where there are arcs $x\to y$ with label $k$ and
$y\to x$ with label $l$ as follows:
\begin{itemize}
\item If $k=l=1$, let $x$ and $y$ be distinct elements in an abelian group $G$.
\item If $k=l>1$, let $x=b$, $y=ab$ in the dihedral group of order $2^{k+1}$.
\item If $k=2$ and $l>2$, let $x=a$, $y=b$ in the dihedral group of order
$2^{l+1}$.
\item If $2<k<l$, take $G$ to be the direct product of dihedral groups of orders
$2^{k+1}$ (with generators $a_1$ and $b_1$) and $2^{l+1}$ (with generators
$a_2$ and $b_2$), with $x=b_1a_2$ and $y=a_1b_1b_2$.
\end{itemize}

Now suppose that $k$ is finite and greater than $1$, while $l$ is infinite.
Take $G_1$ to be dihedral of order $2^{k+1}$, and $G_2$ to be dihedral of 
order $6$ (with $c$ of order~$3$ and $d$ of order $2$). In $G_1\times G_2$,
take $x=bc$ and $y=abd$, noting that $(c,d)$ has infinite Engel depth.

Finally, suppose that both $k$ and $l$ are infinite. We can take $G$ to be
the direct product of two copies of the dihedral group of order $6$, with 
$x=c_1d_2$ and $y=d_1c_2$.

\medskip

So the theorem is proved for $2$-vertex subgraphs, and we need to extend to
subgraphs with larger numbers of vertices.

For each pair of vertices $u,v$ of $D$, choose a group realising the induced
sub-digraph on $\{u,v\}$ according to the arguments just given. There is no
overlap between groups for different pairs. Let $G(u,v)$ be the group
corresponding to the vertices $u$ and $v$, and let $g_u$ and $g_v$ be the
group elements which are identified with those vertices. (A little care is
required; in some cases the two elements $x$ and $y$ are symmetric in the
group, we have to specify which is associated with which vertex of $D$.)

Take $G$ to be the direct product of these groups over all pairs $u,v$. Now
take a vertex $u$.  We identify it with an element of the direct product as
follows: in any factor indexed by a pair containing $u$, we choose the element
$g_u$; in any other factor, we choose the identity. Clearly different vertices
of $D$ are identified with different group elements. Now when we evaluate an
Engel commutator, say $[y,{}_kx]$, on a pair $x=u$, $y=v$, the commutator
evaluates correctly in the group $G(u,v)$; for any other group, at least one
of $x$ and $y$ has been identified with the identity, and so even the first
commutator is trivial. Thus the Engel depth of the group elements chosen to
represent $u$ and $v$ agrees with the number on the arc $u\to v$ if there is
one, and is infinite if there is no such arc.

Thus, the theorem is proved.
\end{proof}

\end{document}
